\chapter{Conclusions}
\label{conclusions}

\par This thesis compared C and Rust by building Gentlix Bank twice: one React frontend, 
one PostgreSQL schema, and two interchangeable backends that implement the same 22 REST 
routes, JWT format, and banking rules. That setup made it possible to study memory management, 
design patterns, safety, tooling, migration effort, and performance on equal product terms 
rather than on toy programs.

\par \textbf{Memory and architecture.} Chapter~\ref{chap:ch3} showed that C gives full 
control at the price of manual allocation, NULL discipline, and lifetime tracking. Rust 
moves many failure modes to compile time through ownership, borrowing, and \texttt{Option}. 
For a ledger application, the relevant question is not only whether the program crashes, 
but whether it can corrupt data silently---a theme developed fully in Chapter~\ref{chap:ch5}.

\par \textbf{Design patterns.} Chapter~\ref{chap:ch4} demonstrated that both languages 
can implement encapsulation, polymorphic repositories, composition, and layered architecture. 
In C these are conventions enforced by headers and review; in Rust they are language rules 
(\texttt{pub}, \texttt{impl}, traits, module boundaries). The architectures match; the 
verification story differs.

\par \textbf{Safety and concurrency.} Chapter~\ref{chap:ch5} showed that Rust reduces manual 
obligations---152 \texttt{free} sites, 198 nullable uses, hand-rolled JWT, manual column 
reads in C versus \texttt{Option}, \texttt{Drop}, crates, and \texttt{sqlx} macros in Rust 
(Table~\ref{tab:audit_obligations})---while preserving the same financial behaviour when 
both backends follow SQL transactions and domain checks. Statement replay and parallel 
analytics illustrate where mutex discipline and thread join order matter for correct dashboard 
numbers, not only for crash-free operation.

\par \textbf{Ecosystem and API.} Chapter~\ref{chap:ch6} documented how \texttt{libpq} column 
indices versus \texttt{sqlx} macros shape schema safety, and how pagination on the statement 
endpoint (\texttt{limit}, \texttt{offset}, \texttt{total\_count}) is part of the shared 
contract both backends must honour.

\par \textbf{Development and runtime performance.} Chapter~\ref{chap:ch8} combined LOC and 
audit tables with build-system comparison and local k6 benchmarks. C led on latency and 
throughput on the tested machine while Rust carried lower manual verification load.

\par \textbf{Future work.} Fair release-mode Rust benchmarks, repeated runs on AWS behind 
the load balancer described in Chapter~\ref{sec:ch2sec4}, and expanded automated tests on 
both backends would strengthen the performance chapter. Further work could also explore 
formal sanitiser runs on the C tree and gradual replacement of individual C modules rather 
than a full dual implementation.

\par In summary, Gentlix Bank supports a practical conclusion: Rust and C can deliver the 
same online banking experience, but they distribute correctness work differently---C toward 
runtime discipline and review, Rust toward compile-time proof and ecosystem support---and 
teams modernising a C backend should weigh both the safety audit savings and the measured 
runtime cost in their own deployment context.
